\textbf{Voraussetzungen}:
\begin{enumerate}
    \item $D \subset \mathbb R^n$ offen
    \item $w: I=[a,b] \rightarrow D$ diffbare Kurve
    \item $f: D \rightarrow \mathbb R$ stetiges Skalarfeld oder \\
          $f: D \rightarrow \mathbb R^n$ stetiges Vektorfeld
\end{enumerate}

\subsection{Gradientenfelder}
\begin{sectionbox}
    \subsubsection{Wegzusammenhängend}
    Eine Menge ist wegzusammenhängend, wenn man von jedem Punkt zu jedem anderen Punkt kommt, ohne die Menge zu verlassen (keine gespaltenen Mengen ohne Verbindungen)
\end{sectionbox}

\begin{sectionbox}
    \subsubsection{Einfach wegzusammenhängend}
    Eine Menge ist einfach wegzusammenhängend, wenn man jede geschlossene Kurve auf einen Punkt zusammenziehen kann, ohne die Menge zu verlassen.
    Hierbei darf man die Kurve beliebig verschieben.
    Ist z.B. eine Linie im $\mathbb R^3$ eine Definitionslücke, so ist die Menge nicht einfach zusammenhängend, da man eine Kurve um die Linie nicht zusammenziehen kann (Linie unendlich hoch)
\end{sectionbox}

\begin{sectionbox}
    \subsubsection{Wichtige Sätze}
    Sei $D\subset \mathbb R^n$ eine offene, \textbf{wegzusammenhängende} Menge.
    Sei $f: D\rightarrow \mathbb R^n$ ein stetiges Vektorfeld.
    Dann ist äquivalent
    \begin{cookbox}{Äquivalente Aussagen}
        \item $f$ ist Gradientenfeld
        \item Es gibt eine Stammfunktion
        \item Für alle stetig diffbaren Kurven $w: [a,b] \rightarrow D$ hängt das Kurvenintegral nur von $w(a)$ und $w(b)$ ab
        \item $\oint_w f \cdot dx=0$ für alle geschlossenen, stetig diffbaren Kurven $w$
        \item (bei $\mathbb R^3$) $\rot f=0$
    \end{cookbox}

    \noindent\makebox[\linewidth]{\rule{\columnwidth}{0.4pt}} \\

    Sei $D \subset \mathbb R^n$ eine offene, \textbf{einfach wegzusammenhängende} Menge.
    Sei $f: D\rightarrow \mathbb R^n$ ein stetiges Vektorfeld.
    Dann gilt
    \begin{emphbox}
        \setlength{\tabcolsep}{0.1em}
        \begin{tabular}{cl}
            f ist Gradientenfeld $\Longleftrightarrow$ & $J_f(x)=J_f(x)^T \; \forall x \in D$ \\
                                                       & $\rot f=0$ in $\mathbb R^3$
        \end{tabular}
    \end{emphbox}

\end{sectionbox}
